% docs/manual/operations/chapters/05-cache.tex
% 第 5 章：缓存策略

\chapter{缓存策略}

缓存是 \openbench{} 演进中"省时间"的关键设计。本章详解三层缓存的协作、多用户共享、清理时机、损坏检测。

\section{三层缓存回顾}

用户卷第 6 章已概念性介绍。运维视角补充：

\begin{longtable}{p{2.5cm} p{4cm} p{6.5cm}}
  \toprule
  层 & 存活时长 & 覆盖范围 \\
  \midrule
  \endhead
  Memory & 单次 \cli{run} 进程内 & 本进程访问的 NC 文件 + 中间数组 \\
  Disk (zarr) & \file{output/<run>/data/} 持久 & 预处理后对齐的 sim/ref 数据 \\
  Task cache & \file{output/<run>/scratch/cache/} & 任务级 hash → 输出文件清单 \\
  \bottomrule
\end{longtable}

三层独立失效：memory 跨 run 不保留；disk zarr 跨 run 命中（如配置一致）；task cache 跨 run 命中（如配置一致），但在 cache 丢但 zarr 还在时仅省 metric 计算。

\section{Cache key 算法}

\file{runner/cache.py:make\_cache\_key} + \texttt{EvaluationCache.hash\_config}：

\begin{minted}{python}
import hashlib
import json

def hash_config(d: dict) -> str:
    canonical = json.dumps(d, sort_keys=True, default=str)
    return hashlib.sha256(canonical.encode()).hexdigest()
\end{minted}

\textbf{覆盖 cache 的字段}（变 → invalidate）：

\begin{itemize}
  \item \texttt{variable}：变量名
  \item \texttt{sim\_source}：sim 标签
  \item \texttt{ref\_source}：ref 名
  \item \texttt{metrics}：指标列表
  \item \texttt{scores}：打分列表
  \item \texttt{comparisons}：comparison item 列表
  \item \texttt{statistics}：statistic item 列表
\end{itemize}

\textbf{不覆盖 cache 的字段}：

\begin{itemize}
  \item CLI 临时 \cli{--cores}（不影响输出值）
  \item CLI 临时 \cli{--variables}（缩 task 列表，不改单 task 行为）
  \item \yamlkey{num\_cores}（同上）
  \item \yamlkey{debug\_mode}
\end{itemize}

\warninline{\textbf{潜在 bug}：如果 \openbench{} 升级，metric 算法可能变；但 cache key 不含 \openbench{} 版本，旧 cache 会被误命中产出旧版数值。\emph{每次升级建议手动清 cache 或加 \yamlkey{force: true}}。这是 v3.0 待优化点。}

\section{多用户共享 cache}

HPC 上同一个 reference 数据集可能被多个用户/项目反复评估。共享 cache 能省大量时间。

\subsection{共享方案 A：共享 zarr 但独立 task cache}

\begin{minted}{yaml}
project:
  output_dir: /scratch/shared/openbench/eval-${USER}-${DATE}
\end{minted}

每个用户独立输出目录，但底层 zarr 数据可以共享一个目录（通过软链接或共享 \texttt{data/}）。

\subsection{共享方案 B：完全独立但 reference 数据指向共享路径}

\openbench{} 默认每次 run 在自己目录写 zarr。要让多次 run 共享，把 reference 数据集本身放在共享只读区，让 \texttt{registry.root\_dir} 指向共享路径。\openbench{} 直接读 NC，不重复写 zarr（如果 zarr 也是共享的）。

\subsection{并发写冲突}

\openbench{} 的 cache 写是 \emph{atomic}（先写 \texttt{.tmp}，rename 替换）。两个用户同时写同名 cache 文件时：

\begin{itemize}
  \item 文件级原子：要么 A 的版本，要么 B 的，不会损坏
  \item 不阻塞：先到先得，后者覆盖前者
  \item \emph{结果}：如果 A 与 B 配置一致，cache 内容一致，无问题；不一致则后写者胜
\end{itemize}

\section{cache\_cleanup.py：自动清理}

\file{util/cache\_cleanup.py}（216 行）提供清理工具。

\subsection{LRU + 大小限制}

模块支持基于最近访问时间的清理：

\begin{minted}{python}
from openbench.util.cache_cleanup import cleanup_cache

cleanup_cache(
    cache_dir="output/eval-2026/scratch",
    max_size_gb=50,      # 总大小限额
    max_age_days=30,     # 最大保留天数
)
\end{minted}

实际策略：

\begin{itemize}
  \item 删除超过 \texttt{max\_age\_days} 的 task cache
  \item 然后按 LRU 删除直到总大小 ≤ \texttt{max\_size\_gb}
  \item 写日志记录每次清理
\end{itemize}

\subsection{手动触发}

当前 v3.0a1 没有 \cli{openbench cache clean} 命令；用 Python 脚本：

\begin{minted}{python}
# clean_old_caches.py
from openbench.util.cache_cleanup import cleanup_cache
from pathlib import Path

for run_dir in Path("/scratch/shared/openbench").iterdir():
    if (run_dir / "scratch").exists():
        cleanup_cache(
            str(run_dir / "scratch"),
            max_size_gb=20,
            max_age_days=14,
        )
\end{minted}

\section{何时该清缓存}

\begin{longtable}{p{5cm} p{8cm}}
  \toprule
  情形 & 清什么 \\
  \midrule
  \endhead
  \openbench{} 升级 & 全部 cache（metric 实现可能变） \\
  Reference catalog 改 & 受影响 reference 的 zarr + task cache \\
  发现某次 run 数值异常 & 该 run 的 \texttt{scratch/} \\
  磁盘配额超限 & 旧的 run 整个目录 \\
  异常断电 / Ctrl-C 频繁 & 怀疑的 task 的 cache 文件 \\
  从 partial 状态重跑 & 失败 task 的 cache（让它重新计算） \\
  \bottomrule
\end{longtable}

\section{zarr 格式细节}

\openbench{} 内部用 \modname{xarray.to\_zarr}。每个 task 的 zarr 包含：

\begin{itemize}
  \item \texttt{sim}：simulation 数据 array
  \item \texttt{ref}：reference 数据 array
  \item 元信息（坐标、单位、标识）
\end{itemize}

\subsection{Chunking}

默认 chunk 由 xarray/dask 决定。grid 评估典型 chunk：\texttt{(time=12, lat=180, lon=360)}。改 chunk 影响内存与读速度。

\subsection{压缩}

zarr 默认 blosc 压缩。压缩比通常 5-10x。

\subsection{跨版本兼容}

\openbench{} 内部 zarr schema 可能在大版本间变化。\textbf{不保证跨大版本兼容}。升级后清旧 zarr 是安全做法。

\section{cache 损坏检测与修复}

\subsection{症状}

\begin{itemize}
  \item 某 task 反复失败但其他 task 正常
  \item zarr 读取时报 \texttt{KeyError} 或 \texttt{InvalidArrayError}
  \item 中断后 zarr 写入半截
\end{itemize}

\subsection{诊断}

\begin{minted}{bash}
# 看 zarr 元数据
python -c "
import zarr
g = zarr.open('output/<run>/data/<task>.zarr', mode='r')
print(g.tree())
print(list(g.array_keys()))
"
\end{minted}

如果报错，cache 损坏。

\subsection{修复}

\begin{minted}{bash}
# 删损坏的 zarr
rm -rf output/<run>/data/<task>.zarr

# 删对应的 task cache
rm output/<run>/scratch/cache/<task_key>.json

# 重跑：会重新生成
openbench run config.yaml
\end{minted}

\openbench{} 的 cache 检测当前 \emph{依赖损坏时报错}，没有主动 checksum 验证。如果担心静默损坏，定期 \texttt{rm -rf output/<run>/data} 强制重建。

\section{下一步}

下一章详解大规模站点评估实战：千站点资源预算、HDF5 lock 历史与现状、失败站点诊断。
